

\lussec 4. Chiral symmetry relations

Since the quark flow equations (2.4) preserve chiral symmetry,
the time-dependent composite fields like the densities (2.8) 
transform in a simple way under chiral rotations
and thus allow the spontaneous breaking of the symmetry
to be studied in useful new ways.
In this section, the principal goal is to relate
the time-dependent condensates (2.17) to the 
ordinary chiral condensate and the pseudo-scalar decay
constant in the chiral limit.


\lussubsec 4.1 Quark field equation and PCAC relation

The space-time dimension is now set to 4
and the need for regularization and renormalization
is ignored for the moment, i.e.~the statements made 
in this subsection
are, strictly speaking, only valid
at tree level of perturbation theory. 

Perhaps somewhat unexpectedly, the quark field equation
in the theory in 4+1 dimensions,
\equation{
  \langle\{(\slash{D}+M_0)\psi(x)-\lambda(0,x)\}\kern1pt
  \phi_1(t_1,x_1)\ldots\phi_n(t_n,x_n)
  \rangle=
  \hbox{contact terms},
  \enum
}
differs from the one
in QCD by a term that couples to the quark fields at non-zero
flow time. Equation (4.1) holds for any
local fields $\phi_1,\ldots,\phi_n$ and the contact terms vanish
if all points $(t_1,x_1),\ldots,(t_n,x_n)$ in 4+1 dimensions 
are different from $(0,x)$.
The validity of the field equation (4.1) follows
from the identities
\equation{
  (\slash{D}+M_0)\wick{\psi(x)\psibar(y)}
  =\longwick{\lambda(0,x)\psibar(y)}+\delta(x-y),
  \enum
  \nexteq{2.5ex}
  (\slash{D}+M_0)\wick{\psi(x)\chibar(s,y)}
  =\longwick{\lambda(0,x)\chibar(s,y)},
  \enum
  \nexteq{2.5ex}
  (\slash{D}+M_0)\wick{\psi(x)\lambdabar(s,y)}
  =\longwick{\lambda(0,x)\lambdabar(s,y)},
  \enum
}
which are a direct consequence of the formulae quoted
in subsect.~3.3 for the Wick contractions of the fermion fields.
Equations (4.2)--(4.4) moreover show that contact terms only arise 
from contractions of the fundamental quark fields.

Now let 
\equation{
  \Aax^{rs}_{\mu}(x)=\psibar_r(x)\dirac{\mu}\dirac{5}\psi_s(x),
  \qquad
  \Pax^{rs}(x)=\psibar_r(x)\dirac{5}\psi_s(x),
  \enum
}
be the axial currents and densities in QCD.
For $r\neq s$, the quark field equation then implies the
PCAC relation
\equation{
  \langle
  \{\partial_{\mu}\Aax^{rs}_{\mu}(x)-(m_{0,r}+m_{0,s})\Pax^{rs}(x)
  +\Ptax^{rs}(x)\}\kern1pt
  \phi_1(t_1,x_1)\ldots\phi_n(t_n,x_n)\rangle
  \noenum
  \nexteq{2.0ex}
  \quad=\hbox{contact terms},
  \enum
}
where $m_{0,r},m_{0,s}$ are the bare masses of the quarks 
with flavour labels $r,s$ and
\equation{
  \Ptax^{rs}(x)=\lambdabar_r(0,x)\dirac{5}\psi_s(x)
  +\psibar_r(x)\dirac{5}\lambda_s(0,x).
  \enum
}
In particular, the identity
\equation{
  \langle
  \{\partial_{\mu}\Aax^{rs}_{\mu}(x)-(m_{0,r}+m_{0,s})\Pax^{rs}(x)
  +\Ptax^{rs}(x)\}\kern1pt P^{sr}_t(y)\rangle=0
  \enum
}
holds at all flow times $t>0$ and all points $x$ including $x=y$.
There are no contact terms in this case, because
the fermion propagators that contribute to the 
correlation function on the left of eq.~(4.8) all go
from $(0,x)$ to $(t,y)$ in 4+1 dimensions or from $(t,y)$ to $(0,x)$.
The contraction (4.2) thus does not occur. 
In view of 
the smoothing character of the flow equations, all parts
of the correlation function are in fact regular functions 
of $x$.

Another purely algebraic consequence of the structure of the 
Wick contractions is the relation
\equation{
  \int\rmd^4x\,\langle\Ptax^{rs}(x)P^{sr}_t(y)\rangle
  =\langle S^{rr}_t(y)\rangle+\langle S^{ss}_t(y)\rangle,
  \enum
}
which, when combined with the PCAC relation (4.8),
leads to an identity
\equation{
  (m_{0,r}+m_{0,s})\int\rmd^4x\,\langle\Pax^{rs}(x)P^{sr}_t(y)\rangle
  =-\Sigt^{rr}-\Sigt^{ss}
  \enum
}
that allows the time-dependent quark condensates to be 
related to the physics of the pseudo-scalar mesons
(see subsect.~4.3). In the derivation of this equation,
both the quark masses and the flow time 
were assumed to be positive, as otherwise
the absence of boundary terms at large $x$
and non-integrable singularities at $x=y$ 
would not be guaranteed.


\lussubsec 4.2 Renormalized PCAC relation

Beyond tree level of perturbation theory, the fields and quark masses
in eqs.~(4.8)--(4.10) require renormalization. The equations then become
relations among renormalized correlation functions, whose exact
form depends on the chosen normalization conventions.
There are some natural choices one can make here
and it is the goal in the following lines to specify these.
The discussion implicitly assumes a sensible regularization
of the theory, such as dimensional regularization or 
Wilson's lattice formulation [\cite{Wilson}],
where most symmetries are preserved. As before, only the 
flavour non-diagonal channels are considered.

The
non-singlet axial currents and densities
renormalize multiplicatively,
\equation{
  \rens{\Aax}{\mu}^{rs}=Z_A\Aax_{\mu}^{rs},
  \qquad
  \ren{\Pax}^{rs}=Z_P\Pax^{rs},
  \enum
}
and the same is the case for the time-dependent densities
\equation{
  \rens{S}{t}^{rs}=Z_{\chi}S_t^{rs},
  \qquad
  \rens{P}{t}^{rs}=Z_{\chi}P_t^{rs}.
  \enum
}
As already mentioned, since there are
no short-distance singularities at positive flow time,
these fields (including the flavour diagonal
ones) renormalize according to their quark content, i.e.~like 
the product of a time-dependent quark and antiquark field at non-zero 
distance in 4+1 dimensions.

In the case of the field $\Ptax^{rs}$,
the multiplicative renormalizability,
\equation{
  \ren{\Ptax}^{rs}=\tilde{Z}_P\Ptax^{rs},
  \enum
}
is less obvious,
because the field has dimension $4$ and could mix
with other fields of dimension $3$ and $4$.
The Wick contractions involving the
Lagrange-multiplier fields $\lambda$ and $\lambdabar$ 
are however such that, up to contact terms,
the two-point correlation functions of
$\Ptax^{rs}$ and any local 
field composed from the gauge and quark fields 
at flow time zero vanish. 
Divergent additive renormalizations by such fields are therefore
excluded. One is then left with
fields involving the Lagrange-multiplier fields, but 
since $\Ptax^{rs}$ is the only one among these with
dimension 4 and the required symmetry properties,
additive renormalizations are again not possible
and the field must consequently renormalize multiplicatively.

The normalizations of the renormalized fields and 
the renormalized quark masses $\mR{r}$ may now be chosen such that
the PCAC relation assumes the form
\equation{
  \langle
  \{\partial_{\mu}\rens{\Aax}{\mu}^{rs}(x)-
  (\mR{r}+\mR{s})\ren{\Pax}^{rs}(x)
  +\ren{\Ptax}^{rs}(x)\}\kern1pt\rens{P}{t}^{sr}(y)\rangle=0.
  \enum
}  
Note that this convention only fixes the relative normalization of 
the axial current $\rens{\Aax}{\mu}^{rs}$, the product
$(\mR{r}+\mR{s})\ren{\Pax}^{rs}$ and the density $\ren{\Ptax}^{rs}$.
The normalization of the latter may however
be fixed by requiring the equation
\equation{
  \int\rmd^4x\,\langle\ren{\Ptax}^{rs}(x)\rens{P}{t}^{sr}(y)\rangle
  =\langle\rens{S}{t}^{rr}(y)\rangle+\langle\rens{S}{t}^{ss}(y)\rangle
  \enum
}
and thus the identity
\equation{
  (\mR{r}+\mR{s})\int\rmd^4x\,
  \langle\ren{\Pax}^{rs}(x)\rens{P}{t}^{sr}(y)\rangle
  =-\rens{\Sigma}{t}^{rr}-\rens{\Sigma}{t}^{ss}
  \enum
}
to hold (where $\rens{\Sigma}{t}^{rr}=Z_{\chi}\Sigma_t^{rr}$).
Further normalization conditions then still need to be imposed
on the renormalized 
quark masses and the time-dependent densities, but for the masses
one can adopt one of the standard conventions, while the normalization of
the time-dependent fields tends to cancel in the equations of interest 
and is left unspecified.

The normalization of the axial currents
implicitly determined by eqs.~(4.14),(4.15) coincides with
the one usually chosen, where the axial charges
assume integer values. There are different ways to 
show this, the perhaps most straightforward one starting from
a formulation of lattice QCD that preserves chiral symmetry 
[\cite{GinspargWilson}--\cite{ExactChSy}].
Using chiral Ward identities [\cite{BoulderReview}],
the bare fields can be shown to satisfy eqs.~(4.8),(4.9) in this case,
up to terms vanishing proportionally to the lattice spacing.
Equations (4.14),(4.15) then imply
$Z_A=\tilde{Z}_P=1$ and thus the canonical normalization of 
the renormalized axial currents. In view of the universality of the
continuum limit,
the canonical normalization is therefore guaranteed 
by these equations independently 
of the chosen regularization of the theory.


\lussubsec 4.3 Pseudo-scalar matrix elements and the chiral condensate

The quark multiplet is now assumed to contain two light quarks, 
referred to as the $u$ and $d$ quarks, with the same
renormalized mass $\ren{m}$.
If there are not too many further
quarks, chiral symmetry in the $(u,d)$-channel 
is expected to be spontaneously broken and there is a triplet of 
pseudo-scalar mesons, the ``pions", whose mass vanishes
as $\ren{m}$ goes to zero.

At large distances, the light-quark pseudo-scalar correlation
functions are dominated by the intermediate one-pion states.
In particular, as $x_0\to\infty$
\equation{
  \int\rmd^3{\mib x}\,\langle\ren{\Pax}^{ud}(x)\ren{\Pax}^{du}(0)\rangle
  =-{\Gpi^2\over\Mpi}
  \rme^{-\Mpi x_0}\bigl\{1+\rmO(\rme^{-\Delta Ex_0})\bigr\},
  \enum
  \nexteq{2.5ex}
  \int\rmd^3{\mib x}\,\langle\rens{\Aax}{0}^{ud}(x)\ren{\Pax}^{du}(0)\rangle
  =\Fpi\Gpi\rme^{-\Mpi x_0}\bigl\{1+\rmO(\rme^{-\Delta Ex_0})\bigr\},
  \enum
}
where $\Delta E$ denotes the energy gap to the higher states and
$\Mpi$, $\Fpi$ and $\Gpi$ are the pion mass, 
decay constant and vacuum-to-pion matrix element of the axial 
density.

From the point of view of the theory in 4 dimensions,
$\rens{\Pax}{t}^{du}(x)$ is not a local field, but
it is still a well localized expression in the fundamental fields
since the smoothing kernel $K(t,x;0,y)$ falls off approximately 
like a Gaussian at large distances $|x-y|$.
At times $x_0$ much larger than the smoothing range $\sqrt{8t}$, 
the two-point function
\equation{
  \int\rmd^3{\mib x}\,\langle\ren{\Pax}^{ud}(x)\rens{\Pax}{t}^{du}(0)\rangle
  =-{\Gpi\Gpit\over\Mpi}
  \rme^{-\Mpi x_0}\bigl\{1+\rmO(\rme^{-\Delta Ex_0})\bigr\}
  \enum
}
therefore decays in the same way as the ordinary pseudo-scalar
correlation functions.
The coefficient $\Gpit$ coincides with
the vacuum-to-pion matrix element of a certain operator,
but is here considered to be defined through eq.~(4.19).

The low-energy effective constants characterizing the 
(two-flavour) chiral limit,
\equation{
  \Sigma=\lim_{m_{\rm R}\to0}\Fpi\Gpi,
  \enum
  \nexteq{2.5ex}
  F=\lim_{m_{\rm R}\to0}\Fpi,
  \enum
}
are the chiral condensate $\Sigma$
and the value $F$ of the pion decay constant at $\ren{m}=0$.
As already mentioned, these can be related to 
the time-dependent light-quark con\-den\-sate
$\rens{\Sigma}{t}=\rens{\Sigma}{t}^{uu}$ via
the chiral Ward identity (4.16). Using the PCAC
relation
\equation{
  2\ren{m}\Gpi=\Mpi^2\Fpi,
  \enum
}
the latter may be written in the form
\equation{
  \rens{\Sigma}{t}=-{\Mpi^2\Fpi\over2\Gpi}
  \int\rmd^4x\,\langle \ren{P}^{ud}(x)\rens{P}{t}^{du}(0)\rangle.
  \enum
}
The integral on the right of this equation 
diverges when $\ren{m}\to0$ and its asymptotic behaviour is determined
by the pion pole of the correlation function in momentum space.
Recalling eq.~(4.19), one is then led to conclude that
\equation{
  \Sigma=\lim_{m_{\rm R}\to0}{\rens{\Sigma}{t}\Gpi\over\Gpit},
  \enum
  \nexteq{2.5ex}
  F=\lim_{m_{\rm R}\to0}{\rens{\Sigma}{t}\over\Gpit}.
  \enum
}
In the chiral limit, the time-dependent condensate $\rens{\Sigma}{t}$
thus provides an order parameter for the spontaneous
breaking of chiral symmetry, which coincides with the standard
condensate $\Sigma$ up to a
proportionality constant.

A remarkable feature of eqs.~(4.24),(4.25) is the fact that
there is no reference to the axial currents, not even implicitly 
via the 
normalization conditions, and that the equations hold for 
any (positive) flow time $t$.
Moreover, the renormalization constant 
$Z_{\chi}$ drops out so that the decay constant 
$F$, for example, is directly given through the bare time-dependent 
condensate and pseudo-scalar vacuum-to-pion matrix element.


\lussubsec 4.4 Chiral perturbation theory

Close to the chiral limit, the asymptotic behaviour of 
many quantities in QCD can be described by chiral perturbation theory.
If the Compton wavelength of the pion is much larger than
the smoothing range of the gradient flow, i.e.~if
\equation{
  8t\Mpi^2\ll1,
  \enum
}
the time-dependent densities $\rens{S}{t}^{rs}$ and $\rens{P}{t}^{rs}$
are local fields from the point of view of pion physics, with
the same symmetry properties as the densities at vanishing flow time.
At low energies, the correlation functions of the time-dependent 
fields and the ordinary axial currents and densities are
therefore expected to be dominated by the 
pion intermediate states and should hence be accessible to a quantitative
description in the framework of chiral perturbation theory.

The time-dependent densities can in fact easily be included in 
chiral perturbation theory by adding the appropriate source terms 
to the chiral Lagrangian [\cite{ChPT}]. At each order of the 
chiral expansion, these terms require a number of effective couplings
to be introduced, which will, in general, depend on the flow time.
The expansion of the chiral condensate $\rens{\Sigma}{t}$
and the matrix element $\Gpit$, for example, 
can then be worked out
and provides insight into how exactly the
chiral limits (4.24),(4.25) are reached.
